% Draft LaTeX manuscript for the final crop-5% multitask experiment.
% NOTE: This is a clean starting point (not MDPI-formatted yet).

\documentclass[11pt]{article}

\usepackage[margin=1in]{geometry}
\usepackage{graphicx}
\usepackage{booktabs}
\usepackage{amsmath}
\usepackage{siunitx}
\usepackage{hyperref}

\title{Joint prediction of infection state and progression time from microscopy using a multi-task ResNet with temporal reliability profiling}
\author{(Authors omitted for draft)}
\date{\today}

\begin{document}
\maketitle

\begin{abstract}
Automated analysis of microscopy time courses can support infection phenotyping by predicting both infection state and disease progression stage. We propose a multi-task deep learning approach that jointly predicts infection status and post-infection time from microscopy images. To further assess reliability over time, we compute performance within sliding temporal windows.

Because the data provider (GMU) noted an approximate \textasciitilde5\% field-of-view overlap between adjacent acquisition positions, we adopt a simple, robust leakage-mitigation preprocessing step in the final experiment: we remove the outer 5\% border of each frame (center crop; retaining the central 90\%\,$\times$\,90\% region) prior to resizing.
\end{abstract}

\section{Introduction}
(Write intro: HBMVEC infection microscopy, need state+stage prediction, limitations of single-task models.)

\section{Methods}
\subsection{Task definition}
We train a multi-task model to predict (1) infection state (infected vs. mock/PBS) and (2) progression time (hours) from microscopy frames.

\subsection{Leakage mitigation via border cropping (final setting)}
To reduce potential near-duplicate content leakage caused by the reported \textasciitilde5\% field-of-view overlap, we center-crop each frame to remove the outer 5\% border before resizing to the model input size. This preprocessing is applied consistently to train/validation/test.

\subsection{Model and training}
(Insert: ResNet50 backbone, 2 heads, loss functions, optimizer details, augmentations, CV protocol.)

\section{Results}
\subsection{Overall cross-validation performance}
Table~\ref{tab:overall_metrics} summarizes the final (crop-5\%) 5-fold cross-validation performance.

% Auto-generated from paper_table_overall_metrics.csv
% Overall metrics for the final (crop-5%) multitask experiment.

\begin{table}[t]
\centering
\caption{Overall 5-fold cross-validation performance for the proposed multi-task model (final crop-5\% preprocessing). Reported as mean $\pm$ standard deviation across folds, along with min/max.}
\label{tab:overall_metrics}
\begin{tabular}{lrrrr}
\hline
Metric & Mean & Std & Min & Max \\
\hline
Classification accuracy & 0.99253 & 0.00096 & 0.99117 & 0.99389 \\
Classification precision & 0.99945 & 0.00111 & 0.99723 & 1.00000 \\
Classification recall & 0.98544 & 0.00283 & 0.98214 & 0.99038 \\
Classification F1 & 0.99239 & 0.00099 & 0.99099 & 0.99380 \\
Classification AUC & 0.99983 & 0.00005 & 0.99975 & 0.99992 \\
Regression MAE (h) & 1.17622 & 0.09701 & 1.04282 & 1.29194 \\
Regression RMSE (h) & 1.49116 & 0.10206 & 1.35687 & 1.62636 \\
\hline
\end{tabular}
\end{table}


\subsection{Temporal reliability profiling}
We evaluate classification performance within sliding temporal windows (window size 6\,h, stride 3\,h). Table~\ref{tab:temporal_metrics} reports mean $\pm$ standard deviation across folds.

% Auto-generated from paper_table_temporal_metrics.csv
% Temporal window classification metrics for the final (crop-5%) multitask experiment.

\begin{table}[t]
\centering
\caption{Temporal reliability profiling using sliding windows (window size $6$ h, stride $3$ h). Each row reports mean $\pm$ standard deviation across 5 folds for the classification task within a given window center time.}
\label{tab:temporal_metrics}
\begin{tabular}{rccccc}
\hline
Window center (h) & Accuracy & F1 & AUC & Precision & Recall \\
\hline
3  & 0.9443 $\pm$ 0.0043 & 0.9350 $\pm$ 0.0057 & 0.9908 $\pm$ 0.0026 & 0.9947 $\pm$ 0.0107 & 0.8825 $\pm$ 0.0170 \\
6  & 0.9646 $\pm$ 0.0051 & 0.9634 $\pm$ 0.0057 & 0.9964 $\pm$ 0.0014 & 0.9958 $\pm$ 0.0084 & 0.9333 $\pm$ 0.0182 \\
9  & 0.9938 $\pm$ 0.0039 & 0.9937 $\pm$ 0.0039 & 1.0000 $\pm$ 0.0000 & 1.0000 $\pm$ 0.0000 & 0.9875 $\pm$ 0.0078 \\
12 & 0.9990 $\pm$ 0.0021 & 0.9990 $\pm$ 0.0021 & 1.0000 $\pm$ 0.0000 & 1.0000 $\pm$ 0.0000 & 0.9979 $\pm$ 0.0042 \\
15 & 1.0000 $\pm$ 0.0000 & 1.0000 $\pm$ 0.0000 & 1.0000 $\pm$ 0.0000 & 1.0000 $\pm$ 0.0000 & 1.0000 $\pm$ 0.0000 \\
18 & 1.0000 $\pm$ 0.0000 & 1.0000 $\pm$ 0.0000 & 1.0000 $\pm$ 0.0000 & 1.0000 $\pm$ 0.0000 & 1.0000 $\pm$ 0.0000 \\
21 & 1.0000 $\pm$ 0.0000 & 1.0000 $\pm$ 0.0000 & 1.0000 $\pm$ 0.0000 & 1.0000 $\pm$ 0.0000 & 1.0000 $\pm$ 0.0000 \\
24 & 1.0000 $\pm$ 0.0000 & 1.0000 $\pm$ 0.0000 & 1.0000 $\pm$ 0.0000 & 1.0000 $\pm$ 0.0000 & 1.0000 $\pm$ 0.0000 \\
27 & 1.0000 $\pm$ 0.0000 & 1.0000 $\pm$ 0.0000 & 1.0000 $\pm$ 0.0000 & 1.0000 $\pm$ 0.0000 & 1.0000 $\pm$ 0.0000 \\
30 & 1.0000 $\pm$ 0.0000 & 1.0000 $\pm$ 0.0000 & 1.0000 $\pm$ 0.0000 & 1.0000 $\pm$ 0.0000 & 1.0000 $\pm$ 0.0000 \\
33 & 1.0000 $\pm$ 0.0000 & 1.0000 $\pm$ 0.0000 & 1.0000 $\pm$ 0.0000 & 1.0000 $\pm$ 0.0000 & 1.0000 $\pm$ 0.0000 \\
36 & 1.0000 $\pm$ 0.0000 & 1.0000 $\pm$ 0.0000 & 1.0000 $\pm$ 0.0000 & 1.0000 $\pm$ 0.0000 & 1.0000 $\pm$ 0.0000 \\
39 & 1.0000 $\pm$ 0.0000 & 1.0000 $\pm$ 0.0000 & 1.0000 $\pm$ 0.0000 & 1.0000 $\pm$ 0.0000 & 1.0000 $\pm$ 0.0000 \\
42 & 1.0000 $\pm$ 0.0000 & 1.0000 $\pm$ 0.0000 & 1.0000 $\pm$ 0.0000 & 1.0000 $\pm$ 0.0000 & 1.0000 $\pm$ 0.0000 \\
45 & 1.0000 $\pm$ 0.0000 & 1.0000 $\pm$ 0.0000 & 1.0000 $\pm$ 0.0000 & 1.0000 $\pm$ 0.0000 & 1.0000 $\pm$ 0.0000 \\
\hline
\end{tabular}
\end{table}


\subsection{Figures}
(Insert your final figures here. For example:)

\begin{figure}[t]
\centering
\includegraphics[width=0.95\linewidth]{classification_by_time_window.png}
\caption{Classification performance by temporal window (window size 6\,h, stride 3\,h).}
\label{fig:cls_by_window}
\end{figure}

\begin{figure}[t]
\centering
\includegraphics[width=0.95\linewidth]{prediction_scatter.png}
\caption{Regression scatter plot: predicted vs. true progression time.}
\label{fig:reg_scatter}
\end{figure}

\section{Discussion}
(Discuss: low early-stage recall, saturation after 9h, implications, remaining leakage risks.)

\section{Conclusion}
(One paragraph.)

\end{document}
